% chapters/ch_part_xiii_world_simulation_vision.tex
% Part XIII: World Simulation Vision — The Ideal Form
%
% This chapter documents what the world simulation OUGHT to be when fully
% realised. Not what currently runs (Stage 1 prototype: 100k agents, single
% civilisation, 5 world types). The Vision is a 1-billion-agent multi-world
% multi-civilisation simulation with full individual layer, cultural
% pluralism, theory pluralism, and counterfactual mode.
%
% Author note: This chapter is a *forward-looking specification document*.
% Implementation is incremental (Stage 1 -> Stage 5 = v∞). Including this
% Vision in the monograph preserves the design intent across the gap
% between Stage 1 reality and full realisation.

\part{World Simulation Vision (Forward-Looking Specification)}

\noindent\textit{Source: Direct authorial specification from Zarame
(NRMO research lead). This Part documents the \emph{ideal form} of the
NRMO ecosystem, of which the Stage 1 prototype (Part XI Chapter
\ref{ch:historical-validation} and the v4.1 Lineage Monte Carlo) is
the first running realisation. The implementation roadmap from Stage 1
to Stage $\infty$ is described in
Section~\ref{sec:vision-roadmap}.}

\chapter{The World Simulation Vision}
\label{ch:vision}

\nrmobox{Document scope (CRITICAL)}{
This chapter \textbf{specifies} a system not yet implemented at full
scale. The current implementation reaches \textbf{Stage 1} of \textbf{5}
stages: 100,000 agents in a single civilisation across 5 world variants.
The Vision describes \textbf{Stage~$\infty$}: 1 billion agents, multiple
co-evolving civilisations, full individual life layer, pluralistic
decision theories, counterfactual history mode, and unknown-contact
worlds.\\[0.4em]
The Vision is documented here so that:
\begin{itemize}
\item Future development has a target architecture rather than ad-hoc
extension;
\item Readers (including future selves) understand the research
ambition that motivated incremental Stage 1 work;
\item The honest gap between Stage 1 reality and Stage~$\infty$ goal is
visible, preventing accidental over-claim.
\end{itemize}
}

% =====================================================================
\section{Architectural Vision}
\label{sec:vision-architecture}

\subsection{The Telescope Architecture}
\label{sec:vision-telescope}

A 1-billion-agent simulation cannot give every agent identical
computational resolution; not on any plausible 2026 hardware, and
arguably not on any future hardware while preserving the per-agent
fidelity required for "people who decide independently." The
\textbf{Telescope Architecture} resolves this by stratifying agents
into four resolution tiers:

\begin{table}[ht]
\centering
\small
\begin{tabularx}{\linewidth}{XlrXl}
\toprule
\textbf{Tier} & \textbf{Resolution} & \textbf{N} & \textbf{Information}
& \textbf{Use} \\
\midrule
Spotlight & 1 day & 1 person & full event log, full memory, full
decision theory & main subject \\
Local Field & 1 month & $\approx 10^3$ & per-agent state, mid-detail
events & social network \\
Civilisation & 1 year & $\approx 10^7$ & vectorised state arrays,
agent classes & civilisation-level dynamics \\
Global Background & 10 years & $\approx 10^9$ & cohort
representations (100 per cohort) & aggregate global dynamics \\
\bottomrule
\end{tabularx}
\caption[Telescope Architecture tiers]{Four-tier Telescope
Architecture. Total \emph{counted} agents reach $10^9$; full-detail
simulation effort scales as $\sim10^6$.}
\label{tab:vision-telescope}
\end{table}

\paragraph{Spotlight protocol.}
At the start of each simulation, the operator selects a Spotlight
target (one individual, identified by birth year, civilisation, family
seed, and birth order). The simulation runs at maximum resolution for
this individual; reduced resolution for everyone else. The Spotlight
can be re-bound mid-simulation, allowing alternative perspectives on
the same simulated world.

\paragraph{Why this matters for NRMO.}
NRMO's central claim is that governance under absorbing failure risk
generalises across scales (individual, family, civilisation). The
Telescope Architecture is the minimal mechanism that lets all three
scales be simulated \emph{simultaneously and mutually consistently} ---
the Spotlight individual's choices feed into family dynamics, which
feed into civilisation aggregates, which loop back as world-level
shocks affecting the Spotlight individual.

\subsection{Multi-World Generator}
\label{sec:vision-multiworld}

Real history is a single sample from a much larger space of possible
worlds. The Vision implements a parametric generator over alternative
worlds:

\begin{table}[ht]
\centering
\small
\begin{tabularx}{\linewidth}{Xl}
\toprule
\textbf{World variant} & \textbf{Characterising parameters} \\
\midrule
Normal Earth & Historical baseline (calibrated to recorded shocks). \\
Science-Heavy & Tech inflection 168 years earlier; religion strength 0.2. \\
Religion-Heavy & Tech inflection deferred; religion strength 0.9; faith
buffer 0.20. \\
Sci-Religion-Mix & Both science and religion strong (1.2 / 0.7). \\
Accelerated-Tech & Industrial revolution at 1500 (300 yr early); tech
acceleration 2.0. \\
Unknown-Encounter & External contact event in 1480--1520 window;
catalyst tech transfer. \\
Custom & User-supplied parameter vector. \\
\bottomrule
\end{tabularx}
\caption[World variants]{World variants. Each variant is a parameter
override on the Normal Earth baseline.}
\label{tab:vision-worlds}
\end{table}

\paragraph{Why multi-world.}
A single-world simulation cannot tell us which features of NRMO's
performance are due to the structure of NRMO and which are artefacts
of a particular historical sequence. Running NRMO across the world
parameter space identifies:

\begin{enumerate}
\item Which world parameters break NRMO (showing its boundary of
applicability);
\item Which competing decision theories dominate in worlds different
from ours (Faith dominates in Religion-Heavy worlds; NRMO dominates
in Science-Heavy and Mix worlds);
\item Whether \textbf{NRMO advantage} is robust or brittle across
worlds.
\end{enumerate}

This is the honest framing required by the
\texttt{honest claim regimen} that has guided this research from the
beginning. Stage 1 has confirmed that NRMO is not universally
dominant; it is dominant in worlds resembling our own and in
science-leaning futures, but Faith outperforms NRMO in Religion-Heavy
and Accelerated-Tech worlds.

\subsection{Cultural Modules}
\label{sec:vision-cultural-modules}

Each civilisation in the simulation is represented by a
\textbf{Cultural Module} that encapsulates region-specific
institutions:

\begin{itemize}[itemsep=2pt]
\item \texttt{Japan\_Module}: honke/bunke/yōshi-saki (multi-branch
portfolio), jisha-azuke, za, buke-hōkō, Edo occupational matrix.
\item \texttt{China\_Module}: zōngzú (extended kinship), kējǔ
(imperial examination), dézhì (rule by virtue), wàiqī (consort
families).
\item \texttt{Europe\_Module}: primogeniture (eldest son inherits),
guild apprenticeship, monasteries, feudal vassalage, princely
states.
\item \texttt{Islamic\_Module}: waqf (charitable endowments),
ulama scholar networks, merchant trade networks, family-based
craft inheritance.
\item \texttt{SubSaharan\_Module}: extended kinship
(\textit{patronymic}), age-grade systems, oral tradition transmission.
\item \texttt{Indic\_Module}: caste-based occupational
heritability, joint-family households, dharma-grounded duty
ethics.
\item \texttt{Indigenous\_Americas\_Module}, \texttt{Polynesian\_Module},
\texttt{Steppe\_Nomadic\_Module}, etc.
\end{itemize}

\paragraph{Module interface.}
Each Cultural Module exposes a common interface:

\begin{lstlisting}[language=Python,basicstyle=\ttfamily\small,frame=single]
class CulturalModule:
    occupations: list[Occupation]       # local job categories
    inheritance_rules: InheritanceFn    # who inherits what
    region_matrix: ndarray              # geographic mobility
    social_classes: list[Class]         # vertical strata
    marriage_rules: MarriageFn          # who marries whom
    shock_response: ShockFn             # cultural response to crisis
    decision_theory_dist: dict          # how strategies are distributed
\end{lstlisting}

This permits multiple modules to coexist (for inter-civilisation
interaction) while preserving each culture's idiosyncratic Special.

\paragraph{Why this matters.}
The current Stage 1 implementation hard-codes Japan-specific
parameters. The Vision generalises this so the simulation runs across
civilisations \emph{simultaneously}, and the operator can ask:
``How does NRMO perform if applied universally vs.\ if applied only in
the civilisation that invents it?''

% =====================================================================
\section{The Individual Layer}
\label{sec:vision-individual}

The current Stage 1 implementation tracks individuals only as
state-vector summaries (occupation, region, six-dim state). The Vision
treats individuals as first-class entities with rich subjective and
biographical detail.

\subsection{Individual State Schema}
\label{sec:vision-individual-state}

\begin{table}[ht]
\centering
\small
\begin{tabularx}{\linewidth}{lXl}
\toprule
\textbf{Field} & \textbf{Description} & \textbf{Example} \\
\midrule
\texttt{birth\_year} & Birth year in simulation & 1842 \\
\texttt{civilisation} & Cultural module & "Japan" \\
\texttt{region} & Region within civilisation & "Setouchi-Kibi" \\
\texttt{family\_seed} & Reference to family lineage & UUID \\
\texttt{birth\_order} & 1st son, 2nd daughter, etc. & 3 \\
\texttt{health} & Vitality, decreasing with age and shock & 0.62 \\
\texttt{skill\_specialty} & Concrete skill (not just occupation class) &
"silk weaving" \\
\texttt{social\_capital} & Connections in social graph & 0.34 \\
\texttt{psychological\_traits} & Risk tolerance, trust, etc. & vec[0.4,
0.7, 0.2, ...] \\
\texttt{decision\_theory} & The theory they personally use & "NRMO\_vNext" \\
\texttt{life\_event\_log} & Log of major events & list[Event] \\
\texttt{memory} & Past lessons and family wisdom inherited & list[Memory] \\
\texttt{subjective\_experiences} & Emotional records of key moments & list[Experience] \\
\bottomrule
\end{tabularx}
\caption[Individual state schema]{Vision-level individual state.
Stage 1 tracks only fields with vectorisable correspondence
(occupation, region, 6-dim aggregate state); the Vision exposes the
full schema for at least the Spotlight individual.}
\label{tab:vision-individual-state}
\end{table}

\subsection{Decision Moments}
\label{sec:vision-decision-moments}

Real human lives are punctuated by a small number of decisive moments.
Each Spotlight individual encounters approximately:

\begin{itemize}[itemsep=1pt]
\item Marriage choice (timing, partner, region of partner's family).
\item Career transition (continue family business, learn new craft,
emigrate to city, enter clergy).
\item Inheritance acceptance (accept main heir status, decline in
favour of sibling, leave the family).
\item Migration (stay in birthplace, relocate to capital, emigrate to
distant province / colony / overseas).
\item Crisis response (war evacuation, plague flight, famine
adaptation, religious conversion).
\item Children allocation (how many sons enter family business
vs.\ marry into other families vs.\ enter temple/monastery).
\end{itemize}

Each Decision Moment is resolved by the individual's decision theory,
applied to their specific subjective state. This is where NRMO,
EVMax, Faith, etc.\ produce different choices for the \emph{same}
external situation.

\subsection{Subjective Experience Logging}
\label{sec:vision-subjective}

For the Spotlight individual, the simulation maintains a chronicle of
internal experience:

\begin{itemize}
\item Year 1842 (birth): no record.
\item Year 1855 (age 13, first plague exposure): \emph{terror, gratitude
upon survival}.
\item Year 1862 (age 20, marriage): \emph{relief, social anchoring}.
\item Year 1868 (age 26, Meiji Restoration): \emph{disorientation, new
opportunity sensed}.
\item Year 1880 (age 38, business decision): \emph{deliberation,
choice}.
\item Year 1895 (age 53, son's death in war): \emph{grief, reflection on
purpose}.
\item Year 1910 (age 68): \emph{retrospective on a life lived}.
\end{itemize}

This is the literary substrate of "human-life-as-narrative." It is not
required for NRMO performance evaluation, but it is required for the
research project to remain \emph{about people}.

% =====================================================================
\section{Multi-Theory Pluralism}
\label{sec:vision-pluralism}

\subsection{Decision Theory as Plugin}
\label{sec:vision-plugin}

The Vision treats every decision theory as a plugin conforming to a
common interface:

\begin{lstlisting}[language=Python,basicstyle=\ttfamily\small,frame=single]
class DecisionTheory(ABC):
    @abstractmethod
    def choose(self, state, observations, memory) -> Action:
        ...
    @abstractmethod
    def update_after_outcome(self, action, outcome) -> None:
        ...

class NRMOTheory(DecisionTheory): ...        # NRMO_vNext + Omega Full
class EVMaxTheory(DecisionTheory): ...       # expected-value max
class FaithTheory(DecisionTheory): ...       # providence-based
class FamilyTraditionTheory(DecisionTheory): # follow ancestral wisdom
class ImitationTheory(DecisionTheory): ...   # mimic local successful
class DriftTheory(DecisionTheory): ...       # null model (v3.7)
class CharismaTheory(DecisionTheory): ...    # follow charismatic leader
class IdeologyTheory(DecisionTheory): ...    # follow doctrine
\end{lstlisting}

\subsection{Theory Distribution and Memetic Dynamics}
\label{sec:vision-memetic}

The Vision implements two levels of theory pluralism:

\paragraph{Static distribution.}
At simulation start, each agent is assigned a theory by sampling from
a per-civilisation distribution. This is the level Stage 1 implements.

\paragraph{Memetic dynamics (Vision only).}
Theories themselves \emph{evolve}: agents observe nearby successes and
failures, and may switch theories. The dominant theory in a society
shifts over time, and crisis events accelerate switching. Specifically:

\begin{itemize}
\item After a major collapse (lineage failure, civilisation collapse),
neighbours of the failed agent re-evaluate their theory.
\item After a major success (lineage flourishing, professional
elevation), neighbours partially adopt the successful theory.
\item In stable times, theory inheritance is high (children adopt
parents' theory).
\item Black-swan events (Reformation, Enlightenment, Industrial
Revolution) can flip whole populations between theories.
\end{itemize}

This is the mechanism by which a world's dominant decision theory
becomes \emph{endogenous} to the world's history rather than
exogenously fixed.

% =====================================================================
\section{Random Event Catalog}
\label{sec:vision-events}

\subsection{Event Types}
\label{sec:vision-event-types}

The Vision includes a catalog of $\sim$50 randomisable historical
events spanning four categories:

\begin{itemize}
\item \textbf{Natural}: volcanic eruptions (Tambora 1815-style),
earthquakes, tsunamis, droughts, Little Ice Age, climate inflection
points.
\item \textbf{Epidemic}: Black Death, smallpox waves, Spanish flu,
cholera, plague reservoirs.
\item \textbf{Technological}: printing press, gunpowder, steam engine,
electricity, antibiotics, computing, biotechnology.
\item \textbf{Political/Cultural}: religious reformations, dynastic
collapses, ideologies (Enlightenment, nationalism, communism,
neoliberalism), unknown-civilisation contact.
\end{itemize}

\subsection{Frequency Modes}
\label{sec:vision-frequency}

The operator selects one of three frequency modes per simulation
run:

\begin{itemize}
\item \textbf{Sporadic}: Each event type's occurrence probability is
calibrated to historical record (median run: 8--12 events occur).
\item \textbf{Frequent}: Probabilities scaled $\times 3$. Events become
common; civilisation must absorb continuous turbulence.
\item \textbf{Sparse}: Probabilities scaled $\times 0.3$. Few events
occur; civilisation drifts under near-stable conditions.
\end{itemize}

The frequency mode is itself a key world parameter. Civilisations may
be optimised for one frequency regime and fail catastrophically in
another. \textbf{NRMO's drift control mechanism is hypothesised to be
particularly valuable in the Frequent regime}, where many small shocks
sum into long-horizon environmental degradation.

\subsection{Black Swan Events}
\label{sec:vision-black-swan}

Beyond the catalog, the Vision implements a \textbf{Black Swan
generator}: events with very low probability but very high impact, not
in the catalog. These include:

\begin{itemize}
\item Sudden technological discontinuity (faster-than-historical
inventions);
\item Civilisation-ending astrophysical events (asteroid strike,
gamma-ray burst, supernova);
\item Unknown-contact events (alien encounter, parallel-civilisation
contact, etc.);
\item Sudden ideology birth (a new theory emerging that does not exist
in the plugin catalog).
\end{itemize}

These are sampled at probability $\sim 10^{-4}$ per simulation year.
Their inclusion forces the simulation to stress-test decision theories
against truly unforeseen circumstances.

% =====================================================================
\section{Counterfactual History Mode}
\label{sec:vision-counterfactual}

The Vision supports a counterfactual mode in which the operator selects
a specific historical decision point and overrides the actor's
choice:

\begin{itemize}
\item ``Japan does not adopt sakoku (national seclusion) in 1639.''
\item ``Constantinople does not fall in 1453.''
\item ``Industrial Revolution begins in Song Dynasty China (1100).''
\item ``Aztec Empire encounters Europeans first as merchants
(1480) instead of conquerors (1519).''
\item ``Tokugawa Yoshimune fails the Kyōhō Reforms (1716).''
\end{itemize}

For each counterfactual, the simulation runs forward from the override
point and produces a \emph{counterfactual world history}. The operator
can then compare:

\begin{itemize}
\item Civilisation trajectory (R, E, G, O, K, X over time);
\item Spotlight individual's life chronicle in counterfactual world
vs.\ original;
\item Strategy ranking (which decision theories thrive in the
counterfactual);
\item Random event distribution (counterfactual histories may produce
different events).
\end{itemize}

This is the highest-leverage research mode, because it directly
answers ``what does NRMO buy us?'' --- by comparing worlds where NRMO
is and is not adopted at a specific decision point.

% =====================================================================
\section{Implementation Roadmap}
\label{sec:vision-roadmap}

\begin{table}[ht]
\centering
\small
\begin{tabularx}{\linewidth}{ccXc}
\toprule
\textbf{Stage} & \textbf{Agents} & \textbf{Features added} & \textbf{Memory} \\
\midrule
\textbf{Stage 1 (v5.0)} & $10^5$ & 5 worlds, 6 theories, simple
random events, single-civ Spotlight & 4 GB \\
Stage 2 (v5.5) & $10^6$ & + Individual layer expansion, full life
chronicle for 100+ subjects & 16 GB \\
Stage 3 (v6.0) & $10^7$ & + Cultural Modules (Japan + 3 others),
inter-civilisation interaction & 64 GB \\
Stage 4 (v7.0) & $10^8$ & + Memetic dynamics, Black Swan events,
counterfactual mode, GPU acceleration & 128 GB \\
\textbf{Stage 5 (v$\infty$)} & $\boldsymbol{10^9}$ & \textbf{+ Full
Telescope across 10+ Cultural Modules; complete subjective experience
logging for Spotlight; full counterfactual matrix} & 256 GB \\
\bottomrule
\end{tabularx}
\caption[Implementation roadmap]{Five-stage roadmap. Stage 1 complete
(this document); Stages 2-5 pending. Memory and runtime estimates
reflect 2026 hardware extrapolations.}
\label{tab:vision-roadmap}
\end{table}

\paragraph{Status as of v5.5 monograph.}
\textbf{Stage 1 complete}: 100,000-agent simulation across 6 world
variants confirms that decision theory dominance is world-dependent.
NRMO dominates in Science-Heavy, Mix, and Unknown-Encounter worlds;
Faith dominates in Religion-Heavy and Accelerated-Tech worlds; Drift is
universally last. This validates the Vision's core hypothesis that
multi-world testing is necessary for honest assessment.

% =====================================================================
\section{Stage 1 Empirical Result}
\label{sec:vision-stage1-result}

\begin{table}[ht]
\centering
\small
\begin{tabularx}{\linewidth}{XCCCCCCC}
\toprule
\textbf{Theory} & \textbf{Normal} & \textbf{Sci-Heavy} & \textbf{Rel-Heavy}
& \textbf{Sci-Rel-Mix} & \textbf{Accel-Tech} & \textbf{Unknown-Enc} & \textbf{Mean} \\
\midrule
Faith & 1.154 & 1.140 & \textbf{1.169} & 1.160 & \textbf{1.623} & 1.222 & 1.245 \\
Adaptive\_OmegaFull & 1.148 & \textbf{1.258} & 1.153 & 1.182 & 1.564 & 1.249 & 1.259 \\
NRMO\_vNext & 1.153 & 1.256 & 1.146 & \textbf{1.186} & 1.565 & \textbf{1.256} & 1.260 \\
RiskAdjustedUtility & 1.145 & 1.213 & 1.143 & 1.155 & 1.531 & 1.225 & 1.235 \\
ExpectedValueMax & 1.140 & 1.149 & 1.150 & 1.148 & 1.338 & 1.156 & 1.180 \\
Drift & 0.817 & 0.837 & 0.818 & 0.822 & 0.850 & 0.817 & 0.827 \\
\bottomrule
\end{tabularx}
\caption[Stage 1 cross-world result]{Composite score per theory per
world variant (n=100{,}000 per cell). \textbf{Bold}: best theory in
column. NRMO/$\Omega$ Full dominates 3/6 worlds; Faith dominates 2/6;
$\Omega$ Full leads in cross-world mean.}
\label{tab:vision-stage1-result}
\end{table}

\paragraph{Key findings (Stage 1).}

\begin{enumerate}
\item \textbf{World dependence is real.} The best theory varies across
world types. NRMO is not universally dominant.

\item \textbf{NRMO has the best cross-world mean} (1.260 vs.\ Faith
1.245), supporting the claim that NRMO is the best general-purpose
governance theory across diverse worlds.

\item \textbf{Faith dominates Religion-Heavy and Accelerated-Tech.}
Religion-Heavy is intuitive (faith-aligned world). Accelerated-Tech is
surprising and warrants further analysis: rapid technological progress
appears to favour low-decision-cost theories (Faith ranks $\#1$ in
Accelerated-Tech).

\item \textbf{Drift is uniformly last.} Across all 6 worlds, the no-agency
baseline ranks last, confirming the agency dividend identified in
v4.1 (lineage simulation).
\end{enumerate}

\paragraph{What Stage 1 cannot say.}

\begin{itemize}
\item No claim about Vision-only features: cultural modules,
inter-civilisation interaction, memetic dynamics, counterfactual mode.
\item No claim about emergent behaviour at $10^9$ agent scale.
\item No claim about specific historical individuals or events.
\item No claim that the world parameters are calibrated to physical
reality.
\end{itemize}

% =====================================================================
\section{Stage 1 Critical Revision: Two-Sided Faith (v5.0.1)}
\label{sec:vision-stage1-revision}

\nrmobox{Critical fix to Stage 1}{
The initial Stage 1 implementation (v5.0) implemented Faith as a single
theory with a unilateral failure-reduction buffer (the ``Faith Buffer'').
This is structurally a \textbf{God's-Invisible-Hand bias}: Faith agents
got a free survival bonus with no costs. The result that Faith dominated
in Religion-Heavy and Accelerated-Tech worlds was therefore a
\textbf{simulation artefact, not a genuine empirical finding}.\\[0.4em]
v5.0.1 fixes this by decomposing Faith into 6 historically-grounded
sub-theories (Buddhist / Communal / Calvinist / Charismatic / Ascetic /
Militant), each with explicit positive AND negative mechanisms; making
\texttt{religion\_strength} an endogenous variable evolving with shock
history and tech acceleration; and implementing endogenous religious
conflict shocks when Militant Faith share is high.
}

\subsection{Six Faith Sub-Theories}
\label{sec:vision-faith-subtheories}

\begin{table}[ht]
\centering
\small
\begin{tabularx}{\linewidth}{lXX}
\toprule
\textbf{Sub-theory} & \textbf{Positive mechanism} & \textbf{Negative mechanism} \\
\midrule
Buddhist & Psychological resilience (impermanence acceptance), literacy promotion (sutras) & Sengoku suppression of Ikkō-ikki, Meiji haibutsu-kishaku \\
Communal & Strong communal insurance (waqf, charity), kinship support, education promotion & None major (historically robust) \\
Calvinist & Discipline, work ethic, scripture literacy push & Persecution by Catholic regimes (Sengoku/Edo eras) \\
Charismatic & Hope-based resilience for individuals & \textbf{Cult collapse tail events} (mass-death possibility) \\
Ascetic & Personal psychological resilience & \textbf{Vow of celibacy} -- structural dead-end for biological lineage \\
Militant & Warrior-brotherhood mutual aid & Combat mortality, \textbf{generates endogenous religious wars when share is large} \\
\bottomrule
\end{tabularx}
\caption[Faith sub-theories with both effects]{Six Faith sub-theories with
both positive and negative mechanisms. Each sub-theory is calibrated to
historical-sociological evidence rather than asserted by fiat.}
\label{tab:vision-faith-sides}
\end{table}

\subsection{Endogenous Religion Strength}
\label{sec:vision-religion-endogenous}

\texttt{religion\_strength} is no longer a fixed world parameter; it
evolves over time:

\begin{itemize}
\item Major shock (epidemic, war) raises religion\_strength
($+0.03 \times$ shock magnitude). Historical analogue: flagellant
movements following the Black Death (1347--1353).
\item Tech acceleration lowers religion\_strength
($-0.008 \times$ excess factor). Historical analogue: Inglehart--Welzel
secularisation thesis.
\item Militant share above 15\% causes polarisation
($+0.01$ per generation).
\item Modern long-run secularisation: $-0.005$ per generation after
1900.
\end{itemize}

\subsection{Endogenous Religious Conflict}
\label{sec:vision-religious-conflict}

When Militant Faith agents form a sufficient minority in a religious
society, they generate \textbf{shocks affecting all agents (not just
Faith agents)}:

\begin{lstlisting}[language=Python,basicstyle=\ttfamily\small,frame=single]
if militant_share > 0.05 and religion_strength > 0.30:
    P_conflict = 0.12 * militant_share * religion_strength * era_mult
    if conflict_triggered:
        global_shock += 0.04 + 0.10 * militant_share * religion_strength
\end{lstlisting}

This mechanism realises the historical truth that ``\emph{religious
wars produce no winners}.'' In Stage 1 v5.0.1 stress tests, when Militant
share is forced to 33\%, three endogenous conflicts add a cumulative
$+0.218$ shock to the world; Militant Faith itself drops from would-be
\#1 to \#6 in ranking, while non-violent Faith\_Communal becomes \#1
(beneficiary of the conflict-induced demand for stability).

\subsection{v5.0.1 Cross-World Result}
\label{sec:vision-stage1-revised-result}

\begin{table}[ht]
\centering
\scriptsize
\begin{tabularx}{\linewidth}{XCCCCCCC}
\toprule
\textbf{Theory} & \textbf{Normal} & \textbf{Sci-Heavy} & \textbf{Rel-Heavy}
& \textbf{Sci-Rel-Mix} & \textbf{Accel-Tech} & \textbf{Unknown-Enc} & \textbf{Mean} \\
\midrule
Faith\_Communal & 1.16 & 1.26 & 1.16 & 1.20 & \textbf{1.75} & \textbf{1.51} & \textbf{1.34} \\
Faith\_Buddhist & \textbf{1.19} & 1.25 & \textbf{1.17} & 1.20 & 1.61 & 1.31 & 1.29 \\
Adaptive\_OmegaFull & 1.15 & 1.26 & 1.13 & 1.18 & 1.56 & 1.25 & 1.25 \\
NRMO\_vNext & 1.15 & 1.25 & 1.14 & 1.17 & 1.54 & 1.26 & 1.25 \\
RiskAdjustedUtility & 1.15 & 1.21 & 1.14 & 1.16 & 1.49 & 1.23 & 1.23 \\
Faith\_Militant & 1.10 & 1.22 & 1.13 & 1.16 & 1.51 & 1.14 & 1.21 \\
Faith\_Calvinist & 1.15 & 1.17 & 1.16 & 1.17 & 1.34 & 1.19 & 1.20 \\
ExpectedValueMax & 1.14 & 1.15 & 1.14 & 1.15 & 1.22 & 1.15 & 1.16 \\
Faith\_Charismatic & 1.06 & 1.15 & 0.94 & 1.02 & 1.61 & 1.06 & 1.14 \\
Drift & 0.82 & 0.83 & 0.81 & 0.83 & 0.82 & 0.82 & 0.82 \\
\textbf{Faith\_Ascetic} & 0.29 & 0.29 & 0.29 & 0.30 & 0.28 & 0.30 & \textbf{0.29} \\
\bottomrule
\end{tabularx}
\caption[Stage 1 v5.0.1 result]{Cross-world ranking with two-sided
Faith implementation. Faith\_Communal achieves the highest mean, but
this reflects historical reality of communal-based religions
(Confucianism, Islamic waqf, early Christianity), not simulation bias.
Faith\_Ascetic correctly suffers structural collapse due to
celibacy-induced biological dead-end.}
\label{tab:vision-stage1-v501}
\end{table}

\paragraph{Honest interpretation.}
The v5.0.1 result reveals three key insights:

\begin{enumerate}
\item \textbf{Communal Faith is robust because it implements the same
distribution-heavy strategy as NRMO\_vNext.} The two are not really in
competition; they are different verbal framings of similar action
allocations.

\item \textbf{NRMO is third in cross-world mean but uniformly robust.}
Unlike Faith\_Communal which dominates only certain worlds,
NRMO\_vNext / $\Omega$ Full perform consistently well in all 6 worlds
without ever achieving runaway success or catastrophic failure.

\item \textbf{Faith\_Ascetic correctly suffers structural collapse.}
The vow of celibacy in monastic traditions produces no biological
heirs. Stage 1 v5.0.1 reproduces this with continuation rates of
17--22\% across all worlds.

\item \textbf{Faith\_Charismatic is world-dependent.} In high-religion
worlds, cult collapse events trigger more often. This is consistent
with historical patterns (Jonestown, Heaven's Gate, etc.\ require a
broader religiously-receptive cultural milieu).

\item \textbf{Faith\_Militant cannot self-perpetuate.} When Militant
share is high, it generates shocks harming itself. This realises the
``no winners in religious wars'' historical pattern.
\end{enumerate}

The v5.0 result that ``Faith universally dominates'' was a
simulation artefact. The v5.0.1 result that ``Faith is heterogeneous;
Communal Faith and NRMO are converging on similar policies; Ascetic
Faith fails structurally; Militant Faith self-undermines'' is
\textbf{the more honest empirical finding}.

% =====================================================================
% Note: "Connections to Other Parts" section moved to ch_part_xiii_v52_to_v62.tex
%       (after the v5.2--v6.2 development chapters)
